\chapter{Panoramica del Network layer}

\section{Introduzione al Network Layer}
Il compito principale del Network Layer è instradare i pacchetti del livello transport, gestendo il trasferimento e la comunicazione host-to-host.



\subsection{Instradare vs Trasferire}
Sono tra le responsabilità principali del livello di rete, e \underline{non sono la stessa cosa}. Non sono sinonimi, ma solo \textbf{concetti collegati}.

\begin{itemize}
	\item \textbf{Routing, instradamento}.\\
	Definiamo come \textbf{routing} il processo di \textbf{selezione} e \textbf{definizione} di \textbf{cammini} (path, route) per i pacchetti all'interno o tra reti. Indica anche i processi relativi alla \textbf{gestione del traffico} generale \textbf{della rete}. \textit{Routing is the process of selecting and defining paths for IP-packet traffic within or between networks as well as the process of managing network traffic overall.}
	
	\item \textbf{Forwarding, trasferimento}.\\
	Avviene dopo il routing, \textbf{consiste nell'inoltro effettivo dei pacchetti}, secondo i termini stabiliti dal routing.
\end{itemize}

\subsection{Data plane vs Control plane}
La vastità degli argomenti che riguardano il Network Layer, ci spingono a suddividere la discussione tra \textbf{data plane} e \textbf{control plane}.

\begin{itemize}
	\item \textbf{Data plane.}\\
	Il dataplane è la parte del router o dello switch che si occupa del trasferimento effettivo dei pacchetti, inoltrandoli da un’interfaccia all’altra secondo i termini stabiliti del control plane. Opera in modo veloce e automatico, spesso in hardware. È quindi il responsabile del forwarding!
	
	\item \textbf{Control plane.}\\
	Il control plane è responsabile di prendere decisioni sul percorso dei pacchetti, costruendo e aggiornando le tabelle di routing seguendo dei protocolli dedicati. Gestisce la logica e la configurazione della rete. È chiaramente il responsabile del routing.
\end{itemize}

\newpage

\section{Routing}

Esistono due approcci differenti:
\begin{itemize}
	\item \textbf{Traditional Routing Algorithms.}\\
	Effettuato dai singoli router, senza una visibilità globale, ma dei propri neighbours. È una struttura distribuita, più resistente della prossima soluzione, ma più lenta. Ogni router ha al proprio interno una tabella d'instradamento, cui valoi vengono aggiornati di volta in volta.
	\item \textbf{Software Defined Network.}\\
	È un'architettura centralizzata esterna, che sa tutto della rete. Ottiene path migliori e più velocemente, ma se l'SDN va giù, il sistema crolla. Inoltre, riceverà un altissimo numero di richieste non-distribuite, in quanto, in questo tipo di architettura, è proprio il controller remoto a calcolare e distribuire le tabelle di inoltro.
	
\end{itemize}
\subsection{Self-Organizing-Routing}
Scambiandosi informazioni, ottengono una visione generale della rete. Determinano un routing migliore tramite algoritmi quali Bellman-Ford e Dijkstra. Questi algoritmi sono rinforzati tramite euristiche, tecniche AI e di reinforcement learning. Andremo ad approfondire questa parte affrontando il control-plane. 

\subsection{SDN - Software Defined Networking}
Il control plane è centralizzato in un remote controller.\\
Conosce tutta la topologia della rete, e prende decisioni sui percorsi che devono seguire i pacchetti, sa come aggiornare le tabelle di inoltro nei router. Manda istruzioni ai router sull'instradamento e riceve aggiornamenti dai router sul traffico. Il control agent (CA) è un modulo che gira sui router che permette di comunicare col Remote Controller. Il remote agent non agisce mai in autonomia dal remote controller.

\subsection{Virtual Circuit Service, circuito virtuale}
È una modalità di commutazione in cui, tra due peer, viene costruito un \textbf{circuito virtuale} in fase di connessione (secondo i soliti metodi di routing). L'inoltro dei pacchetti \textbf{di quella connessione} procederà \underline{esclusivamente} sul \textbf{percorso stabilito}. In queste circostanze è chiaro che i router \textbf{non siano state-less}: manterranno infatti delle informazioni di routing relative a quella connessione. Un qualsiasi \textbf{malfunzionamento} dei router nel percorso, implica una \textbf{riconnessione da zero}. In compenso, il \textbf{controllo della congestione} è molto più \textbf{semplice}. 

\subsection{Packet Switching, commutazione di pacchetto}
La modalità a circuito virtuale è molto diversa alla più usata modalità di \textbf{packet switching} (ovvero della \textbf{commutazione a pacchetto}) in cui tutti i \textbf{pacchetti}, anche della stessa connessione, sono \textbf{inoltrati separatamente} sulla base delle informazioni contenute negli \textbf{header}. Instradano ogni pacchetto secondo le proprie tabelle di routing, e \underline{non pensando ad un circuito fisso tra due host}.


\subsection{Router}
Al giorno d'oggi sono dei "mini-stack ISO/OSI", cui livello più alto è il transport.\\
Attualmente, i router contengono un \textit{firmware}\footnote{Sono programmi scritti nella memoria flash, solitamente in linguaggio di basso livello (Assembly, C di basso livello). Sono software che la CPU può eseguire direttamente grazie al program counter. All'accensione del dispositivo, il codice viene copiato nella SRAM, e la CPU lo esegue. Questi firmware sono velocissimi, ed è il motivo per cui rispondono alle esigenze dei router.} nella propria memoria flash.
\begin{itemize}
	\item Esamina gli \textbf{header} dei datagrammi IP che gli passano attraverso.
	\item Sposta i datagrammi \textbf{dalla porta di input} a quella di \textbf{output} per trasferire i datagramma nel path stabilito. Effettuano un \textbf{forward} dei pacchetti.
\end{itemize}

\subsection{Indirizzamento logico}
Consiste nell'assegnamento di \textbf{indirizzi logici} ai dispositivi (tramite indirizzi \textbf{IP}) per \textbf{identificare} in modo univoco ogni \textbf{nodo nella rete}.

\section{Frammentazione e riassemblaggio}
I \textbf{datagrammi IP troppo grandi} per il router devono essere \textbf{frammentati e riassemblati al mittente} per renderne possibile il trasferimento. Questo è uno dei compiti del data plane del Network Layer (la MTU dell'Ethernet è di 1500 byte). In dei casi, si vorrebbe poter avere la possibilità di controllare se permettere o meno la frammentazione. Questo è ottenuto tramite dei flag opportuni dell'header IPv4 (Don't Fragment e More Fragment). Scopriremo anche che IPv6 non permette la frammentazione da parte dei router intermedi.


\section{Modelli di servizio di rete}
In inglese, \textbf{Network Service Model}.
Stabiliscono le caratteristiche del trasporto dei datagrammi da host a host.\\
I servizi possono fare riferimento ai singoli pacchetti, quali
\begin{itemize}
	\item \textbf{Consegna garantita}.
	\item Consegna garantita \textbf{entro $n$ millisecondi}.
\end{itemize}
O a interi flussi di datagrams, come
\begin{itemize}
	\item \textbf{Consegna} dei datagrammi \textbf{in ordine}.
	\item \textbf{Bandwidth minima} garantita.
	\item \textbf{Maximum jitter}, garantisce che la variazione massima tra i ritardi dei pacchetti.
\end{itemize}

Su Internet, garantire questi servizi è semplicemente utopia. In campo militare o finanziario, alcuni di questi servizi devono essere garantiti. Nelle reti degli ATM, viene usato un servizio di \textbf{Available Bit Rate}. Garantisce una bandwidth minima utilizzando la banda di rete disponibile.

\subsection{Internet Service Model - Best Effort}
In Internet, le \textbf{garanzie} in termini \textbf{di reliability} sono affidate al \textbf{livello di trasporto}, tramite \textbf{RDT} e \textbf{meccanismi vari del TCP}. Il network, nel caso di Internet, non si occupa di assicurare i servizi appena esposti, ma di offrire dei meccanismi di instradamento e inoltro dei pacchetti per i livelli superiori.

\subsubsection{Best Effort - Complice della popolarità di Internet?}
La semplicità dei meccanismi dietro ai servizi \textbf{best-effort di Internet}, \textbf{ha indubbiamente permesso} a quest'ultimo di \textbf{espandersi a macchia d'olio}.\\
Con una larghezza di banda sufficientemente ampia, le applicazioni funzionano piuttosto bene la maggior parte delle volte. \textbf{Good enough!}
